% Tapis de Quillen, with Quillen's offprint and preprint (cote 162-5, typed part): archive cote 162-5, site page 35 (PDF page 36).
% typescript. Draft compared: none.
% AI-transcribed by Claude Opus 5.5 (claude-opus-5-5, adaptive thinking, effort high), Message Batches API, batch t162_full (msgbatch_01UAh3zyWPd454AeDkhnmQ33), 2026-11-13; not yet checked by eye.
\scanpage{162-5/35}
\xout{pleinementxfidèlex(respx} une équivalence). Même énoncé si on met dans le
coup un \xout{Anneauxinex} systême local d'anneaux sur X'. \unsure{Enfin|Ehfin}, f induit \xout{une}
\inserted[typed]{\xout{foncteurxplxfidxx(respx} une équivalence)}
\xout{isomorphisme} entre la catégorie des coefficients locaux sur C et celle
des coéff. locaux sur C' \xout{(??xpeut-êtrexcelaxdoitxil}\inserted[typed]{\xout{cettexpropriétéxdoit}}\xout{xêtrexmisedansxles}
\xout{hypothèsesxpréliminairesxsurxfx?)x} Principe de démonstration: on commence
par prouver ce dernier résultat, en notant que \xout{dansxlexca} si un topos est
locqalement hom. trivial, il est loca connexe et loc. simplement connexe,
d'où une bonne théorie du $\hand{\pi}_o$ et du $\hand{\pi}_1$ (qui sont ici discrets), et on est
ra\unsure{\over{n}{m}}ené \xout{pour} à un cas particulier du critère d'homotopisme de Artin et
Mazur, savoir un critère \inserted[typed]{cohomologique} pour qu'un homomorhisme de groupes G\hand{$\longrightarrow$}H (ici \xout{l}des
groupes $\hand{\pi}_1$ de X,X') soit un isomorphisme: il doit induire des isomorphis\xout{mx}
mes sur les H$^o$ et H$^1$ (y inclus dans le cas non commutatif...). On prouve
la pleine fidélité en se ramenant par la suite spectrale aux cas \unsure{des|les} Ext$^i$
de coéfficients locaux. Par la suite spectrale encore, cela résultera du
fait suivant: si X\hand{\unsure{$'$}} est \unsure{localement} hom. trivial, alors la catégorie des
faisceaux abéliens loc. constants est stable par \uline{Ext}$^i$ , et le foncteur
M \hand{$\longmapsto$}M$_X$ de\xout{sxfaisceauxxconstants} (Ab) dans X$_{ab}$ commute aux dits Ext$^i$ .
En fait, X et X' étant loc homp\hand{|}triviaux, les conditions suivantes sur f
seront équivalentes:

a) f est\unsure{b} un homotopisme, i.e. induit pour tout systême local (pas
néc commutatif) sur X'un isomorphisme sur les H$^i$\xout{pourxixxx}.

b) \xout{Mêmexchosexave} f induit une équivalence D$^b_{lc}$ (X')\hand{$\longrightarrow$}D$^b_{lc}$(X) .

\unsure{a\hand{$'$})} f induit une équivalence entre la catégorie des systêmes locaux
abéliens sur X' et X, et des isomorphismes sur les H$^i$ correspondants (donc
on ne prend ici que des coéfficients commutatifs).

\margin{?} J'ignore si on peut dans a) se borner aux systêmes locaux commutatifs,
\xout{ixexxdans}
L'équivalence entre a) et b) fournit une première justification ou
\mspage{35} motivation pour définir des types d'homotopie via la catégorie D\hand{$^b_{lc}$(X)},
